\section{Сравнение с Self-Collaboration}\label{sec:scc}
Мы сравниваем Self-Collaboration (SCC) с одношаговыми условиями, используя ту же модель Luna, средний уровень рассуждения и штатный оценщик. Это сравнение относится к полным процессам. Адаптированный авторский контроллер~\cite{scc2024source} начинает с плана аналитика в JSON --- текстовом формате структурированных данных, вызывает разработчика для получения Python-кода, а затем просит тестировщика сгенерировать до пяти примеров вызова в функции \texttt{check(candidate)}. Он выполняет программу и сгенерированную проверку; внутренним сигналом остановки служит выполнение без исключений, при этом напечатанные значения автоматически не сравниваются с ожидаемыми ответами. Заданный предел допускает одну первоначальную реализацию и не более одного исправления. Модель получает полный предоставленный контекст задачи и историю сообщений ролей, записанную в текстовом формате. Эталонные решения и тесты бенчмарка остаются внешними.

Первоначальный план содержит 3\,000 блоков «задача--повтор», в каждом из которых назначаются новые процессы SR, SN и SCC. Поправка, принятая после начала генерации, но до проверки качества SCC, фиксирует первые 2\,000 рандомизированных блоков (seed порядка 20260911). Это дает 6\,000 назначений методов на 956 задачах: у 212 задач выбран один повтор, у 444 --- два, у 300 --- три. Из этих задач 941 проходит эталонные проверки. Невыбранные 3\,000 назначений методов не входят в знаменатель после поправки.

Штатное оценивание прошли все 5\,595 завершенных программ: 1\,997 SR, 1\,998 SN и 1\,600 SCC. В сравнении сохранены 397 инфраструктурных сбоев и восемь ранее приостановленных попыток. В частности, у SCC имеется 396 сбоев Docker --- среды запуска изолированных контейнеров --- и четыре приостановленных процесса; у SR --- три приостановленные попытки; у SN --- одна приостановленная попытка и один сбой потока. Ни одна отправленная попытка не генерировалась повторно. После эталонных проверок качество наблюдается для 1\,964 повторов SR, 1\,965 повторов SN и 1\,574 повторов SCC; число успехов равно соответственно 1\,067, 1\,064 и 809. Это количества наблюдаемых повторов, а не независимых задач и не показатели выбора лучшего из трех.

Для каждого сравниваемого метода исходы SCC минус исходы сравниваемого метода сначала усредняются по совпадающим выбранным повторам, в которых наблюдаются оба исхода, а затем с равными весами по допустимым задачам, имеющим хотя бы одну такую пару. Поэтому задача с тремя выбранными повторами имеет тот же вес, что и задача с одним повтором. Двусторонние $t$-тесты на уровне задач для SCC--SR и SCC--SN образуют отдельное семейство из двух тестов с поправкой Холма. Интервалы используют 10\,000 выборок целых задач, генератор PCG64, seed 20260911 и линейные процентильные границы.

Таблица~\ref{tab:scc-quality} показывает парные разности на задачах с наблюдаемыми исходами. Отрицательный знак означает меньшую долю успеха SCC относительно указанного одношагового процесса.
\begin{table}[H]
\caption{Качество SCC на наблюдаемых сопоставленных повторах. Разности и интервалы указаны в процентных пунктах; $n$ --- число участвующих задач.}
\label{tab:scc-quality}
\centering\small
\begin{tabular}{@{}lrrrr@{}}
\toprule
Контраст & $n$ & $\Delta$ & 95\% ДИ & $p$ Холма\\
\midrule
SCC--SR & 883 & $-2{,}68$ & $[-4{,}72;\,-0{,}60]$ & 0,0217\\
SCC--SN & 882 & $-2{,}48$ & $[-4{,}50;\,-0{,}43]$ & 0,0217\\
\bottomrule
\end{tabular}
\end{table}

Оба условных контраста отрицательны (табл.~\ref{tab:scc-quality}). Интервалы чувствительности по группам исходных задач сходны: $[-4{,}72;\,-0{,}58]$ пункта относительно SR и $[-4{,}51;\,-0{,}43]$ относительно SN для объединенного разбиения. Для оценки всей назначенной совокупности отдельно учтены пропуски. Их крайние допустимые значения задают границы разности: каждому неизвестному исходу приписывается сначала наиболее благоприятное, затем наименее благоприятное для SCC значение. Это расчет границ, а не восстановление фактических ответов. Экстремальные границы с фиксированным знаменателем по всем 941 допустимым выбранным задачам равны $[-12{,}98;\,+6{,}61]$ и $[-13{,}14;\,+6{,}43]$ пункта. Первая граница означает, что относительно SR данные допускают как проигрыш SCC до 12,98 процентного пункта, так и выигрыш до 6,61 пункта. Вторая аналогично относится к SN. Таким образом, вся совокупность охвачена интервальной оценкой, а наблюдаемые пары --- более точной условной оценкой. Границы не являются доверительными интервалами: они отражают неизвестные исходы, а не случайный разброс наблюдаемой выборки. Однозначное ранжирование всей совокупности потребовало бы доступных итоговых результатов незавершенных процессов либо дополнительных предположений о пропусках.

Для сравнения ресурсов необходимы оба полностью завершенных процесса с корректными счетчиками. Отношение парных средних по задачам для API-эквивалентной стоимостной оценки равно 3,138 относительно SR (896 задач; 95\%-й интервал $[3{,}067;\,3{,}208]$) и 2,939 относительно SN (895 задач; $[2{,}871;\,3{,}010]$). Эти интервалы отношений средних используют 10\,000 выборок целых задач с seed 20260914. Они относятся к наблюдаемым полным процессам и не объединяют сохранение качества с денежным преимуществом. Доступность и сопоставленные эффекты до и после поправки сохранены как описательные проверки; они не идентифицируют причинный временной эффект.

Отдельная диагностика оценивает все 399 сохраненных программ разработчика из незавершенных процессов SCC, не перезапуская контроллер и не добавляя сгенерированные им проверки. Среди 392 программ, допустимых по контролям, она находит 196 успешных; семь программ не проходят контроль допустимости. Эти результаты промежуточных программ не заменяют неизвестные исходы полного процесса в сравнении.
